\documentclass[10pt, letterpaper]{article}
\usepackage[margin=0.85in]{geometry}
\usepackage{graphicx}
\usepackage{booktabs}
\usepackage{amsmath, amssymb}
\usepackage{hyperref}
\usepackage{xcolor}
\usepackage{caption}
\usepackage{subcaption}
\usepackage{float}
\usepackage{helvet}
\usepackage{enumitem}
\usepackage{tcolorbox}
\usepackage{multirow}
\usepackage{array}
\usepackage{longtable}
\renewcommand{\familydefault}{\sfdefault}

\definecolor{sigcolor}{RGB}{215, 38, 38}
\definecolor{sigbg}{RGB}{255, 235, 235}
\definecolor{nosigcolor}{RGB}{120, 120, 120}
\definecolor{accentblue}{RGB}{0, 114, 189}
\definecolor{keybg}{RGB}{240, 248, 255}

\hypersetup{
    colorlinks=true,
    linkcolor=accentblue,
    urlcolor=accentblue
}

\captionsetup{font=small, labelfont=bf, skip=4pt}

\title{%
    \vspace{-1cm}
    {\Large SEEG Population Dynamics Analysis}\\[0.5em]
    {\large Condition-Dependent Neural Dynamics Report}\\[0.3em]
    {\normalsize Subject: \texttt{%%SUBJECT_ID%%}}
}
\author{Automated Analysis Pipeline --- SEEGring v1.0}
\date{%%DATE%%}

\begin{document}
\maketitle

\begin{abstract}
This report presents a comprehensive analysis of condition-dependent neural dynamics
using stereoelectroencephalography (SEEG) recordings from subject \texttt{%%SUBJECT_ID%%}.
The analysis spans \textbf{%%N_TASKS%% cognitive tasks} (%%N_RUNS%% total runs)
recorded with %%N_CHANNELS%% SEEG channels at %%FS%%~Hz.
All analyses are designed to identify and characterize \textcolor{sigcolor}{\textbf{where and when
neural population dynamics differ across experimental conditions}}.
Statistical significance is assessed via cluster-based permutation tests
($\alpha$ = %%ALPHA%%, %%N_PERMS%% permutations) and reported throughout.
\end{abstract}

\begin{tcolorbox}[colback=sigbg, colframe=sigcolor, title={\textbf{Key Finding: Condition-Dependent Significance}}]
%%KEY_FINDING_BOX%%
\end{tcolorbox}

\tableofcontents
\newpage

%% ======================================================================
\section{Data Summary}
%% ======================================================================

\begin{table}[H]
\centering
\caption{Dataset overview for subject \texttt{%%SUBJECT_ID%%}.
Rows in \textcolor{sigcolor}{red} indicate tasks with statistically significant
condition separation ($p < %%ALPHA%%$, cluster-corrected).}
\begin{tabular}{@{}llrrl@{}}
\toprule
Session & Task & Trials & Conditions & Status \\
\midrule
%%TASK_TABLE_ROWS%%
\bottomrule
\end{tabular}
\label{tab:overview}
\end{table}

%% ======================================================================
\section{Methods}
%% ======================================================================

\subsection{Signal Preprocessing}
Signals were preprocessed with: (1) %%HIGHPASS%%--%%LOWPASS%%~Hz bandpass filtering
(Butterworth, order 4); (2) line noise removal at %%LINE_NOISE%%~Hz and harmonics
via notch filtering; (3) bipolar re-referencing within electrode shafts to maximize
local field potential specificity and minimize common-mode noise.
Bad channels identified via BIDS \texttt{channels.tsv} metadata were excluded
prior to analysis.

\subsection{Epoch Extraction \& Quality Control}
Epochs were time-locked to stimulus onset (%%EPOCH_PRE%%~s pre to %%EPOCH_POST%%~s post).
Trials received traffic-light quality labels (green/yellow/red) based on composite metrics
(RMS amplitude, variance, kurtosis, peak-to-peak). Only green and yellow trials
contributed to analyses, with quality-based weighting.

\subsection{Time-Frequency Decomposition}
Time-frequency representations were computed using Morlet wavelets
(%%TF_CYCLES%% mean cycles, linearly scaled with frequency).
Power was baseline-normalized (%%BASELINE_METHOD%% method) relative to the
pre-stimulus window (%%EPOCH_PRE%%~s to 0~s).

\subsection{Population Dynamics \& Latent Space}
Population activity tensors (trials $\times$ channels $\times$ time) were
z-scored and reduced to %%N_LATENT_DIMS%% latent dimensions via PCA.
Condition-averaged trajectories were smoothed with a %%SMOOTH_MS%%~ms Gaussian kernel.
Trajectory geometry (path length, curvature, tangent angles) and dynamical systems
properties (jPCA rotational dynamics) characterized temporal evolution.

\subsection{Condition Discrimination}
\textbf{Condition-dependent changes} were assessed through multiple complementary approaches:
\begin{enumerate}[nosep]
    \item \textbf{Cluster-based permutation tests} (%%N_PERMS%% permutations, $\alpha$ = %%ALPHA%%)
          on ERP, time-frequency, and latent space measures
    \item \textbf{Time-resolved condition separation index}: Mahalanobis-like distance between
          condition centroids in latent space, with significance assessed via permutation
    \item \textbf{Neural decoding}: LDA classifier with stratified %%N_FOLDS%%-fold
          cross-validation, providing time-resolved classification accuracy
    \item \textbf{Bootstrap confidence intervals}: 500 bootstrap resamples for
          trajectory uncertainty quantification; non-overlapping CIs indicate condition divergence
\end{enumerate}

\noindent
\textcolor{sigcolor}{\textbf{Throughout this report, significant condition-dependent effects
are highlighted}} with pink background shading in temporal plots, contour overlays
in time-frequency maps, bold/colored entries in tables, and dedicated significance panels.

%% ======================================================================
\section{Results by Task}
%% ======================================================================

%%PER_TASK_SECTIONS%%

%% ======================================================================
\section{Condition-Dependent Significance Summary}
%% ======================================================================

\begin{tcolorbox}[colback=keybg, colframe=accentblue, title={\textbf{Cross-Task Condition Discrimination}}]
%%CROSS_TASK_TEXT%%
\end{tcolorbox}

\begin{table}[H]
\centering
\caption{Condition discrimination metrics across all tasks.
\textcolor{sigcolor}{\textbf{Bold red}} entries indicate statistically significant
condition separation ($p < %%ALPHA%%$, cluster-corrected).
Stars: ${}^{*}p<0.05$, ${}^{**}p<0.01$, ${}^{***}p<0.001$.}
\begin{tabular}{@{}l r r r r r r@{}}
\toprule
Task & Peak Sep. & $p_{\min}$ & Peak Acc. (\%) & Onset (ms) & Eff. Dim. & jPCA $R^2$ \\
\midrule
%%CROSS_TASK_TABLE_ROWS%%
\bottomrule
\end{tabular}
\label{tab:significance}
\end{table}

%%SIGNIFICANCE_TIMELINE%%

%% ======================================================================
\section{Significance Interpretation}
%% ======================================================================

%%INTERPRETATION_TEXT%%

\subsection{Methodological Notes}
\begin{itemize}[nosep]
    \item All p-values are corrected for multiple comparisons using cluster-based
          permutation testing (%%N_PERMS%% permutations)
    \item Significance thresholds: ${}^{*}p<0.05$, ${}^{**}p<0.01$, ${}^{***}p<0.001$
    \item Decoding onset defined as first time point exceeding chance level at $p < 0.05$
    \item Condition separation index computed in %%N_LATENT_DIMS%%-dimensional latent space
    \item Effective dimensionality estimated via participation ratio of eigenvalue spectrum
\end{itemize}

\section*{Acknowledgements}
\small
Analysis performed using the SEEGring population dynamics toolbox.
Figures generated with Nature-journal-quality visualization system
(Wong 2011 colorblind-safe palette, Helvetica 5--7pt, 600 DPI).

\end{document}
